\section{The sharp midpoint modulus of the first energy}
\label{sec:r11-sharp-midpoint-modulus}

Retain exactly the strict-box, background-extension and
$\mathcal E^1_{S,\partial}(V)$ hypotheses of the global
first-energy calculus.  Put
\[
 c_q=\frac{1-2^{-q}}{2(q+1)},\qquad
 \mathscr D_S(u,v)=
       \mathscr A_S\left(\frac{u+v}{2}\right)
            -\frac{\mathscr A_S(u)+\mathscr A_S(v)}2 .
 \tag{SM1}
\]
The simpler positive constant in (EV18) already suffices
for the Cauchy argument there.  This section computes
the optimal uniform coefficient and its equality example.

\begin{lemma}[A finite Bernstein coefficient comparison]
\label{lem:r11-midpoint-Bernstein}
For every integer $q\geq1$ and $0\leq l\leq q-1$,
\[
 a_{q,l}:=
 \frac q2\binom{q-1}{l}
    \int_0^1\min(t,1-t)t^l(1-t)^{q-1-l}\,dt
       \geq c_q.
 \tag{SM2}
\]
Equality holds for $l=0$ and $l=q-1$.
\end{lemma}

\noindent\textbf{Status: \texttt{NEW\_PROVED}}
auxiliary finite polynomial inequality; the following
proof holds for every $q$, independently of any numerical
checks.

\begin{proof}
For $1\leq k\leq q$ write
\[
\begin{split}
 p_i(s)&=\binom qi s^i(1-s)^{q-i},\\
 B_k(s)&=\sum_{i=k}^q p_i(s),\\
 b_k(t)&=q\binom{q-1}{k-1}t^{k-1}(1-t)^{q-k}.
\end{split}
 \tag{SM3}
\]
Termwise differentiation of the finite sum telescopes
to $B_k'=b_k$.  The integral of $b_k$ over $[0,1]$
is one.  Also
$B_k(1-s)=1-B_{q-k+1}(s)$, by the binomial theorem.
Consequently
\[
 P_k(s):=\int_s^{1-s}b_k(t)\,dt
              =1-B_k(s)-B_{q-k+1}(s)
                    \quad(0\leq s\leq1/2).
 \tag{SM4}
\]
For $1\leq k\leq(q+1)/2$ this gives
\[
 P_k(s)-P_1(s)
       =\sum_{i=1}^{k-1}\bigl(p_i(s)-p_{q-i}(s)\bigr)\geq0.
 \tag{SM5}
\]
Indeed $i\leq(q-1)/2$, and for $0<s\leq1/2$
the ratio is
$p_i(s)/p_{q-i}(s)=(s/(1-s))^{2i-q}\geq1$.
At $s=0$ use continuity.
The symmetry $P_k=P_{q+1-k}$ proves the same
bound for the remaining $k$.

Layer cake for the nonnegative tent function now gives
\[
 a_{q,k-1}
   =\frac12\int_0^{1/2}P_k(s)\,ds
   \geq\frac12\int_0^{1/2}P_1(s)\,ds=a_{q,0}.
 \tag{SM6}
\]
The endpoint integral, or its reflected version, is
\[
 \frac q2\left[
      \int_0^{1/2}t^q\,dt
        +\int_{1/2}^1(1-t)t^{q-1}\,dt\right]
       =\frac{1-2^{-q}}{2(q+1)}.
 \tag{SM7}
\]
This proves every assertion, including $q=1$.
\end{proof}

\begin{theorem}[Optimal midpoint control of the relative discrepancy]
\label{thm:r11-sharp-midpoint-modulus}
For every $u,v\in\mathcal E^1_{S,\partial}(V)$,
\[
 \boxed{\qquad
       \mathscr D_S(u,v)\geq
          \frac{1-2^{-q}}{2(q+1)}\,\mathscr I_S(u,v).
       \qquad}
 \tag{SM8}
\]
If the class contains a potential of positive first
energy, the constant cannot be increased for that
same domain and background.
No strict positivity or geometric representation of
the original middle-positive $S$ is required.
\end{theorem}

\noindent\textbf{Status: \texttt{NEW\_PROVED}}
as the sharp quantitative form of the global
arbitrary-background midpoint estimate.

\begin{proof}
Let $h=u-v$ and set
\[
 \gamma_l=
       \int_V\Gamma_{N(u^{[l]},v^{[q-1-l]})}(h)
                      \quad(0\leq l\leq q-1).
 \tag{SM9}
\]
Each $\gamma_l$ is finite and nonnegative by (GE7),
and (GE18) gives
$\mathscr I_S(u,v)=\sum_l\gamma_l$.
The scalar Green formula for a twice differentiable
function on $[0,1]$, applied to
$\mathscr A_S((1-t)u+tv)$ and (GE16), yields
\[
 \mathscr D_S(u,v)=
       \frac q2\int_0^1\min(t,1-t)
          \int_V\Gamma_{N(((1-t)u+tv)^{[q-1]})}(h)\,dt
       =\sum_{l=0}^{q-1}a_{q,l}\gamma_l .
 \tag{SM10}
\]
The possible reversal of $l$ under $t\mapsto1-t$
does not affect the coefficient, since the tent
function is symmetric.
All operations are finite sums of the already
proved globally integrable mixed energies.
Apply (SM2) term by term to prove (SM8).

For sharpness choose any $u$ with $e_1(u)>0$ and
take $v=0$.  The zero potential is in the same
boundary power class.  Homogeneity and the definition
of the discrepancy give
\[
 \mathscr I_S(u,0)=e_1(u),\qquad
 \mathscr D_S(u,0)
    =\left(\frac12-2^{-(q+1)}\right)
                       \frac{e_1(u)}{q+1}
    =c_q e_1(u).
 \tag{SM11}
\]
Thus equality occurs for the same original background
and domain, proving optimality.
\end{proof}

\begin{corollary}[An explicit error bound for dual maximization]
\label{cor:r11-sharp-maximizing-error}
In the setting of (EV16), let
$L=\sup\mathcal F_\mu$.
For any $u,v\in\mathcal E^1$,
\[
 \mathscr I_S(u,v)
       \leq \frac{2(q+1)}{1-2^{-q}}\,
                \left[L-\frac{\mathcal F_\mu(u)+\mathcal F_\mu(v)}2\right].
 \tag{SM12}
\]
Thus prescribed objective gaps give a quantitative
bound in every local outer capacity through (EV8).
This remains a convergence estimate; it supplies
neither an ambient classical representative of
a completion limit nor the missing admissible
variations for an unrestricted source equation.
\end{corollary}

\noindent\textbf{Status: \texttt{DIRECT\_COROLLARY}}.

\begin{proof}
The datum integral is affine along the segment,
so it cancels from the midpoint gap.
The value at the midpoint is at most $L$.
Use (SM8) and divide by its positive constant.
\end{proof}

\paragraph{Auxiliary verification.}
The exact rational script
\path{research_r11/verify_midpoint_coefficients.py}
expands the two polynomial integrals for a finite
range of $q$ and checks their minimum against (SM7).
It tests normalization and signs only.
The all-$q$ argument is (SM3)--(SM7), and the analytic
argument is the full original global energy identity
(SM10).
